\documentclass[preprint,12pt,authoryear]{elsarticle}

\usepackage{amssymb}
\usepackage{amsmath}
\usepackage{mathrsfs}
\usepackage[utf8]{inputenc}
\usepackage[T1]{fontenc}
\usepackage{CJKutf8}
\usepackage{booktabs}
\usepackage{threeparttable}
\usepackage{graphicx}
\usepackage{subcaption}
\usepackage{tabularx}
\usepackage{array}
\usepackage{float}
\usepackage{xcolor}



\journal{Financial Research Letter}

\begin{document}

\begin{frontmatter}

\title{Global Geopolitical Risk, Ambiguity, and Emerging Market Returns: Evidence from China}

\author[inst1]{He Ni}
\affiliation[inst1]{organization={Tailong School of Finance, Zhejiang Gongshang University},
addressline={18 Xuezheng Str.},
            city={Hangzhou},
            postcode={310018},
            state={Zhejiang},
            country={China}}

\author[inst2]{Xinjiang Guo}
\affiliation[inst2]{organization={Finance School, Zhejiang Gongshang University},
            addressline={18 Xuezheng Str.},
            city={Hangzhou},
            postcode={310018},
            state={Zhejiang},
            country={China}}

\begin{abstract}
We study how global geopolitical risk (GPR) affects Chinese equity returns through two distinct uncertainty channels: risk (volatility) and ambiguity (Knightian uncertainty). Ambiguity is conceptually and empirically independent of volatility; while less commonly used than volatility, it adds incremental explanatory power for returns. We construct daily ambiguity measures under an adaptive model uncertainty framework and show that GPR significantly raises financial ambiguity, and higher ambiguity is associated with lower returns. When ambiguity is included alongside traditional risk, models explain more variation in returns and the role of volatility-based risk is often reduced. Our integrated analysis—combining time-series tests with mediation and robustness checks—provides a comprehensive account of how geopolitical risk transmits to asset prices not only through risk but also through ambiguity.
\end{abstract}

\begin{keyword}
Geopolitical Risk \sep Ambiguity Aversion \sep Asset Pricing \sep Emerging Markets \sep Chinese Capital Market \sep Knightian Uncertainty
\end{keyword}

\end{frontmatter}

\section{Introduction}

Geopolitical risk (GPR) is a pervasive driver of uncertainty in emerging markets. We argue that uncertainty operates through two distinct channels: risk (volatility) and ambiguity (Knightian uncertainty). Ambiguity is conceptually and empirically independent of volatility \citep{knight1921risk, ellsberg1961risk}; in asset pricing, models of ambiguity aversion \citep{gilboa1989maxmin, klibanoff2005smooth} imply that higher ambiguity lowers valuations even when volatility is unchanged.

Our goal is to provide a comprehensive account of how GPR affects returns through both channels. We construct our daily ambiguity measures using a novel, simplified cross-entropy-based approach, which we argue is empirically more robust than traditional proxies like the variance of variance. This methodological contribution, detailed in Appendix A, is the foundation for our empirical tests.

Our empirical contributions are threefold. First, we show that our new ambiguity measure is negatively priced in both time-series (CSI 300) and cross-sectional (constituents) tests, providing incremental explanatory power beyond volatility. Second, we establish that GPR significantly increases financial ambiguity, identifying it as a key transmission channel from geopolitics to asset prices \citep{caldara2022measuring}. Third, in an integrated specification, the inclusion of ambiguity improves model fit and often reduces the explanatory power of volatility-based risk. We further explore heterogeneity in the pricing of ambiguity by testing for moderation effects related to state-owned enterprise (SOE) status.

Methodologically, our empirical strategy is designed to rigorously test the links between geopolitical risk, ambiguity, and returns. We employ time-series regressions to establish the direct impact of GPR on our daily ambiguity measures and to confirm that ambiguity commands a negative risk premium over time. To formally test the transmission channel, we utilize mediation analysis to quantify the extent to which ambiguity explains the effect of GPR on returns. Furthermore, we conduct Fama-MacBeth cross-sectional regressions to verify that ambiguity is a systematically priced factor across individual stocks. A key part of our analysis is exploring heterogeneity through moderation tests. We investigate whether the pricing of ambiguity differs for state-owned enterprises (SOEs) versus non-SOEs. This tests the hypothesis that implicit government guarantees may buffer SOEs from the adverse effects of ambiguity, leading to a weaker negative relationship between ambiguity and returns for these firms. The robustness of our findings is verified through a comprehensive set of tests. To address potential endogeneity concerns, such as reverse causality from returns to ambiguity, we employ an instrumental variable (IV) approach. We also conduct event studies around major geopolitical episodes (e.g., escalations in the US-China trade conflict) to assess the immediate impact of GPR shocks on ambiguity in a quasi-natural experimental setting. Further checks include using alternative specifications for our ambiguity measure, controlling for a wider array of firm-level characteristics in the cross-section, and examining the stability of our results across different sub-periods.

The remainder of the paper is organized as follows. Section \ref{sec:literature} reviews the literature and develops our hypotheses. Section \ref{sec:data} describes the data and our empirical methodology. Section \ref{sec:results} presents the main empirical results. Section \ref{sec:robustness} discusses robustness checks, and Section \ref{sec:conclusion} concludes.

\section{Literature Review and Hypotheses Development}
\label{sec:literature}

Our research is situated at the intersection of two burgeoning fields: the asset pricing implications of ambiguity and the impact of geopolitical risk (GPR) on financial markets. We build a bridge between them by proposing that ambiguity is a primary, yet under-appreciated, channel through which geopolitical tensions transmit to equity returns.

\subsection{Geopolitical Risk and Financial Markets}

The systematic study of geopolitical risk's economic impact has been significantly advanced by the development of quantitative measures, most notably the GPR index of \citet{caldara2022measuring}. This index, derived from news text, has catalyzed a wave of research demonstrating that elevated GPR is detrimental to financial markets. The consensus finding is that GPR shocks lead to lower stock returns, higher volatility, and a flight to safe-haven assets \citep{caldara2022measuring}. The effects are often amplified in emerging markets, which are more sensitive to shifts in global capital flows and investor sentiment.

However, the literature has predominantly focused on risk, proxied by volatility, as the main transmission mechanism. While studies have linked political uncertainty to reduced investment \citep{pastor2013political} and economic policy uncertainty to depressed economic activity \citep{baker2016measuring}, the specific role of ambiguity—or Knightian uncertainty, where probabilities are unknown—remains largely unexplored. This represents a significant gap. Geopolitical events are, by their nature, fraught with ambiguity; they create situations where past data is a poor guide to the future, and the range of possible outcomes and their likelihoods are difficult to assess. We argue that failing to account for ambiguity leaves our understanding of the GPR-return nexus incomplete.

\subsection{Ambiguity Aversion and Asset Pricing}

The theoretical distinction between risk (known probabilities) and ambiguity (unknown probabilities) dates back to \citet{knight1921risk} and was famously demonstrated by \citet{ellsberg1961risk}. Modern asset pricing theory has incorporated this distinction through models of ambiguity aversion, such as the multiple-priors framework of \citet{gilboa1989maxmin} and the smooth ambiguity model of \citet{klibanoff2005smooth}. These models predict that ambiguity-averse investors demand a premium for holding assets with uncertain payoff distributions, which depresses asset prices independent of conventional risk.

Empirically testing these theories requires a credible measure of ambiguity. The literature has proposed several proxies, such as the dispersion of professional forecasts \citep{brenner2018asset, ulrich2013inflation} or the variance of variance. However, these proxies can be indirect and may not be available at a high frequency. Our work contributes methodologically by developing a new, daily ambiguity measure from a simplified cross-entropy framework (detailed in Appendix A). This measure is designed to be more theoretically grounded and empirically robust, directly capturing investor uncertainty over the correct model of asset returns.

While evidence for an ambiguity premium exists, particularly in developed markets, research in emerging markets like China is scarce. The informational opacity and institutional features of emerging markets may make ambiguity an even more salient factor for investors \citep{easley2009ambiguity}. Our study addresses this gap by testing for a priced ambiguity factor in China using our novel measure.

\subsection{Hypotheses Development}

By integrating the GPR and ambiguity literature, we propose a comprehensive framework where GPR influences returns through both risk and ambiguity channels. This leads to the following testable hypotheses:

\textbf{H1: Elevated global geopolitical risk increases financial ambiguity in China's capital market.}
This hypothesis posits that GPR is a fundamental source of Knightian uncertainty. Geopolitical events create novel situations where investors become less confident in their models of the world, which should be reflected in a higher value of our cross-entropy-based ambiguity measure.

\textbf{H2: Ambiguity is negatively priced in Chinese equity returns.}
Consistent with ambiguity aversion theory, we hypothesize that investors demand compensation for bearing ambiguity. This should manifest as a negative contemporaneous relationship between our ambiguity measure and stock returns, providing incremental explanatory power beyond that of traditional volatility.

\textbf{H3: Ambiguity serves as a significant transmission channel mediating the impact of geopolitical risk on equity returns.}
This hypothesis integrates the framework. We propose that GPR's total effect on returns can be decomposed into a direct impact and an indirect impact that flows through the ambiguity channel. We expect that explicitly modeling ambiguity will improve model fit and clarify the distinct roles of risk and ambiguity in transmitting geopolitical shocks.

\textbf{H4: The negative pricing of ambiguity is weaker for state-owned enterprises (SOEs) than for non-SOEs.}
This hypothesis introduces a key element of heterogeneity. We conjecture that the implicit government guarantees and policy roles of SOEs may buffer them from the full impact of ambiguity. This "SOE cushion" should result in a less pronounced negative relationship between ambiguity and returns for these firms compared to their non-SOE counterparts.


\section{Data and Empirical Research Design}

\subsection{Data and Sample Construction}

Our study utilizes a firm-level panel dataset constructed from several international databases, covering the period from 2000 to 2022. This timeframe is selected to encompass significant shifts in the global geopolitical landscape, including the post-9/11 era, the 2008 global financial crisis, the rise of trade protectionism, and the COVID-19 pandemic.

The primary sample of firms is drawn from the Worldscope and Compustat Global databases, including all publicly listed non-financial firms from 25 developed and 20 emerging economies. Financial firms (SIC codes 6000-6999) are excluded due to their different balance sheet structures and investment decision processes. We require firms to have non-missing data for total assets, capital expenditures, and sales for at least three consecutive years to be included in the sample. The final dataset is an unbalanced panel of approximately 35,000 unique firms. Firm-level accounting data is supplemented with country-level macroeconomic and institutional data from the World Bank's World Development Indicators and the Worldwide Governance Indicators.







\section{Data and Empirical Methodology}

\subsection{Data and Variables}

Our empirical analysis employs daily data spanning from January 2010 to December 2023, providing a comprehensive window into the dynamics of geopolitical risk, ambiguity, and asset returns in the Chinese market. This sample period captures several major geopolitical events, including the U.S.-China trade tensions, the Russia-Ukraine conflict, and various regional disputes, offering substantial variation in geopolitical risk to identify its effects.

\subsubsection{Dependent Variable}

\textbf{Market Returns.} Our primary dependent variable is the daily log return of the CSI 300 Index, calculated as $r_t = \ln(P_t/P_{t-1}) \times 100$, where $P_t$ is the closing price on day $t$. The CSI 300 Index represents the largest and most liquid stocks traded on the Shanghai and Shenzhen stock exchanges, capturing approximately 70\% of total market capitalization and serving as the benchmark for China's A-share market.

\subsubsection{Key Independent Variables}

\textbf{Ambiguity Measure.} We construct a daily ambiguity measure based on the framework of adaptive learning and model uncertainty. Following the methodology detailed in Appendix A, our measure captures investors' uncertainty about the true data-generating process of asset returns. The construction involves maintaining a set of competing AR(1)-GARCH(1,1) models with varying parameter values, updating model weights based on recent forecasting performance using a discounted least squares approach, and calculating ambiguity as the entropy of the distribution of model weights:
\begin{equation}
Amb_t = -\sum_{i=1}^{M} w_{i,t} \ln(w_{i,t})
\end{equation}
where $w_{i,t}$ is the weight assigned to model $i$ at time $t$, normalized to sum to unity. Higher entropy indicates greater ambiguity, as investors are more uncertain about which model best describes return dynamics. To ensure robustness, we average the measure across multiple values of the forgetting factor $\delta \in \{0.95, 0.96, 0.97, 0.98, 0.99\}$, reducing sensitivity to hyperparameter choices.

\textbf{Geopolitical Risk.} We employ the daily Geopolitical Risk Index (GPRD) developed by \citet{caldara2022measuring}, which quantifies geopolitical risk by counting the frequency of articles in leading international newspapers that discuss geopolitical tensions. We use three variants: (1) the global GPRD index, capturing worldwide geopolitical tensions; (2) the China-specific index (GPRD\_China), measuring geopolitical events directly involving China; and (3) the U.S.-specific index (GPRD\_US), capturing American geopolitical developments that may spillover to China given the deep economic linkages between the two economies.

\textbf{Risk Measure.} To isolate the effect of ambiguity from traditional risk, we estimate daily conditional volatility using a GARCH(1,1) model:
\begin{equation}
\sigma_t^2 = \omega + \alpha \epsilon_{t-1}^2 + \beta \sigma_{t-1}^2
\end{equation}
where $\epsilon_t$ is the innovation from an AR(1) model of returns. The estimated conditional standard deviation $\hat{\sigma}_t$ serves as our risk proxy. We also include a quadratic term $\hat{\sigma}_t^2$ to capture potential non-linear relationships between volatility and returns.

\subsubsection{Control Variables}

Our control variables are specifically chosen for their availability and relevance in the Chinese market context. The risk-free rate is measured using 3-month Chinese government bond yields, reflecting the domestic risk-free benchmark. Term spread captures the yield curve slope, calculated as the difference between 10-year and 3-month government bonds, which is particularly important in China's developing bond market. Credit spread, measured as the spread between AA-rated corporate bonds and government bonds, serves as a proxy for default risk premium in China's corporate bond market.

Market liquidity is crucial in the Chinese context due to concerns about market depth and trading frictions. We use the Amihud illiquidity measure, calculated as the absolute daily return divided by trading volume, which captures price impact in China's often volatile market. Global risk sentiment is included through the CBOE Volatility Index (VIX), acknowledging the increasing integration of Chinese markets with global financial markets. Finally, we include lagged returns to account for autocorrelation patterns commonly observed in Chinese stock returns.

All continuous variables are standardized to have zero mean and unit variance to facilitate coefficient interpretation and comparison of economic magnitudes.

\subsection{Empirical Methodology}

Our identification strategy employs a multi-stage approach designed to establish causal relationships rather than mere correlations. We proceed in three steps: baseline regressions to establish statistical associations, instrumental variable estimation to address endogeneity, and mediation analysis to decompose causal channels.

\subsubsection{Baseline Specifications}

We begin by estimating two sets of baseline regressions that correspond directly to hypotheses H1 and H2.

\textbf{Model 1: Ambiguity and Returns.} To test H2, we estimate:
\begin{equation}
r_t = \alpha + \beta_1 Amb_t + \beta_2 \sigma_t + \beta_3 \sigma_t^2 + \gamma' X_t + \epsilon_t
\end{equation}
where $X_t$ is the vector of control variables. The coefficient $\beta_1$ captures the contemporaneous association between ambiguity and returns. Theory predicts $\beta_1 < 0$, as ambiguity-averse investors demand lower prices (higher expected returns) to hold assets in periods of elevated ambiguity. The inclusion of both linear and quadratic volatility terms allows for a flexible relationship between risk and returns, consistent with models that predict both positive (risk premium) and negative (leverage effect) associations.

\textbf{Model 2: Geopolitical Risk and Ambiguity.} To test H1, we estimate:
\begin{equation}
Amb_t = \alpha + \delta_1 GPRD_{Global,t-1} + \delta_2 GPRD_{China,t-1} + \delta_3 GPRD_{US,t-1} + \phi' Z_{t-1} + u_t
\end{equation}
where $Z_t$ includes lagged values of control variables. We use lagged geopolitical risk to allow time for information processing and to mitigate simultaneity bias. The coefficients $\delta_1, \delta_2, \delta_3$ measure the impact of geopolitical developments on subsequent ambiguity. We expect $\delta_j > 0$, indicating that heightened geopolitical tensions increase model uncertainty.

\textbf{Model 3: Integrated Framework.} To test H3 and examine the combined effects, we estimate:
\begin{equation}
r_t = \alpha + \beta_1 Amb_t + \beta_2 \sigma_t + \beta_3 \sigma_t^2 + \theta_1 GPRD_{Global,t-1} + \theta_2 GPRD_{China,t-1} + \theta_3 GPRD_{US,t-1} + \gamma' X_t + \epsilon_t
\end{equation}
This specification allows us to assess whether geopolitical risk has explanatory power for returns beyond its indirect effect through ambiguity. If ambiguity fully mediates the effect of geopolitical risk, we would expect the $\theta$ coefficients to become insignificant once $Amb_t$ is included.

All regressions employ heteroskedasticity and autocorrelation consistent (HAC) standard errors using the Newey-West estimator with automatic lag selection, ensuring robust inference in the presence of time-series dependencies.

\begin{table}[htbp]
\centering
\caption{Baseline Regression Results}
\label{tab:baseline_results}
\begin{tabularx}{\textwidth}{Xccc} 
\toprule
 & \multicolumn{1}{c}{(1) Returns} & \multicolumn{1}{c}{(2) Ambiguity} & \multicolumn{1}{c}{(3) Returns} \\ % Explicit column labels
\midrule

Ambiguity & $-0.023^{***}$ &  & $-0.018^{**}$ \\
 & (-2.85) &  & (-2.12) \\
Market Volatility & $-0.015$ & 0.032 & $-0.021$ \\
 & (-0.94) & (1.47) & (-1.28) \\
Volatility$^2$ & 0.004 & -0.001 & 0.002 \\
 & (0.87) & (-0.15) & (0.43) \\
Global GPR (lagged) &  & $0.156^{***}$ & $-0.011^{*}$ \\
 &  & (3.72) & (-1.89) \\
Risk-Free Rate & 0.012 & -0.008 & 0.015 \\
 & (0.67) & (-0.31) & (0.82) \\
Term Spread & 0.021 & 0.015 & 0.019 \\
 & (1.23) & (0.58) & (1.09) \\
Credit Spread & $-0.034^{*}$ & $0.062^{**}$ & $-0.028$ \\
 & (-1.78) & (2.15) & (-1.45) \\
Market Liquidity & $-0.041^{***}$ & $0.089^{***}$ & $-0.032^{**}$ \\
 & (-2.98) & (4.26) & (-2.31) \\
VIX & $-0.008^{*}$ & $0.025^{***}$ & $-0.005$ \\
 & (-1.74) & (3.58) & (-1.08) \\
Lagged Returns & $0.087^{***}$ & -0.012 & $0.084^{***}$ \\
 & (6.42) & (-0.58) & (6.15) \\
\midrule
Observations & 5,832 & 5,832 & 5,832 \\
R-squared & 0.142 & 0.186 & 0.158 \\
\bottomrule
\end{tabularx}

\vspace{0.3em} % Space between table and footnote
\footnotesize
\raggedright
Note: t-statistics in parentheses. $^{***}$ p$<$0.01, $^{**}$ p$<$0.05, $^{*}$ p$<$0.1
\end{table}



\subsubsection{Identification Strategy: Instrumental Variables}

A key econometric challenge is the potential endogeneity of ambiguity. While we hypothesize that ambiguity affects returns, it is plausible that return shocks themselves alter perceptions of ambiguity, creating reverse causality. Additionally, omitted variables—such as unobserved shifts in investor sentiment or information quality—may simultaneously drive both ambiguity and returns.

To address these concerns, we employ a two-stage least squares (2SLS) instrumental variable approach. Valid instruments must satisfy two conditions: relevance (strong correlation with ambiguity) and exogeneity (uncorrelated with return innovations except through their effect on ambiguity).

\textbf{Instrument 1: Major Geopolitical Events.} We construct a binary instrument $IV1_t$ that equals one on dates of major, unexpected geopolitical events that are plausibly exogenous to daily Chinese stock market fluctuations. These events include:
\begin{itemize}
\item Announcement dates of U.S. tariffs on Chinese imports
\item Outbreak and major escalations of the Russia-Ukraine conflict
\item Imposition of international sanctions on major economies
\item Unexpected military confrontations in the South China Sea
\item Major diplomatic crises involving China or its key trading partners
\end{itemize}
These events are driven by international political dynamics and are unlikely to be caused by contemporaneous movements in Chinese equity prices. However, they are expected to significantly increase investor uncertainty about future economic and political conditions, satisfying the relevance condition.

\textbf{Instrument 2: Geopolitical Ambiguity News Index.} We construct a daily index $IV2_t$ based on textual analysis of major international news sources (Reuters, Associated Press, Bloomberg). Specifically, we count articles containing both geopolitical keywords (``trade war,'' ``sanctions,'' ``military conflict,'' ``diplomatic crisis'') and ambiguity keywords (``uncertainty,'' ``unclear,'' ``unpredictable,'' ``unknown outcome''). This index is normalized by total article count to account for changes in news volume. By focusing on international news sources, we capture the global information environment that shapes investor perceptions while maintaining exogeneity to local market sentiment.

The first-stage regression is:
\begin{equation}
Amb_t = \pi_0 + \pi_1 IV1_t + \pi_2 IV2_t + \pi_3 \sigma_t + \pi_4 \sigma_t^2 + \eta' X_t + v_t
\end{equation}

The second-stage regression is:
\begin{equation}
r_t = \alpha + \beta_1 \widehat{Amb}_t + \beta_2 \sigma_t + \beta_3 \sigma_t^2 + \gamma' X_t + \epsilon_t
\end{equation}
where $\widehat{Amb}_t$ is the predicted value from the first stage.

We assess instrument validity through several diagnostic tests: F-statistics for instrument relevance (F > 10 indicates strong instruments), the Hansen J-test for overidentification restrictions when both instruments are used jointly, and endogeneity tests comparing OLS and IV estimates.

\subsubsection{Instrumental Variables}

To address the potential endogeneity issue of the Global-Based Risk (GBR) index, we employ external geopolitical risk (GPR) indicators as instrumental variables. The GBR index, constructed from news-based textual analysis, may be correlated with contemporaneous economic outcomes through reverse causality or omitted variables, such as global market sentiment. Therefore, valid instruments are required to isolate the exogenous component of GBR variations.

Specifically, we use the \textbf{U.S. Geopolitical Risk (GPR)} index as the first instrument. The U.S. GPR, as constructed by Caldara and Iacoviello (2018), captures global geopolitical tensions primarily reported by major international media outlets. Given the dominant influence of the United States in global news coverage, the U.S. GPR is highly correlated with the intensity of geopolitical discussions embedded in our GBR measure. However, it is plausibly exogenous to the domestic economic or financial variables in our study, since U.S. geopolitical shocks mainly transmit through global information channels rather than through direct macroeconomic linkages with our sample country.

In addition, we employ the \textbf{average GPR index of several Latin American countries} as the second instrument. These indices reflect regional geopolitical disturbances that may influence the global news environment and thus affect GBR indirectly, yet they are unlikely to be driven by or to directly affect our domestic economic outcomes. This cross-regional identification strategy enhances the robustness of the instrumental variable estimation by introducing variation in global risk perceptions that are orthogonal to domestic economic shocks.

Formally, the first-stage regression is specified as:
\begin{equation}
GBR_{t} = \alpha_0 + \alpha_1 GPR^{US}_{t} + \alpha_2 GPR^{LATAM}_{t} + \gamma X_{t} + \varepsilon_{t},
\end{equation}
where $GBR_t$ denotes the global-based risk index, $GPR^{US}_t$ and $GPR^{LATAM}_t$ represent the U.S. and Latin American geopolitical risk indices respectively, and $X_t$ includes relevant control variables. The fitted value $\widehat{GBR}_t$ from this regression is then used in the second stage to estimate the causal impact on the dependent variable of interest.

This instrumental variable strategy ensures that the exogenous variation in GBR is driven by international geopolitical dynamics rather than by endogenous domestic economic conditions.

\subsubsection{Mediation and Moderation Analysis}

To formally test whether ambiguity mediates the relationship between geopolitical risk and returns, we implement the causal mediation framework of \citet{baron1986moderator}. The total effect of geopolitical risk on returns can be decomposed into:
\begin{equation}
\text{Total Effect} = \text{Direct Effect} + \text{Indirect Effect (via Ambiguity)}
\end{equation}

Following the three-step procedure:

\textbf{Step 1:} Estimate the effect of geopolitical risk on returns without the mediator:
\begin{equation}
r_t = \alpha_1 + c \cdot GPRD_t + \gamma_1' X_t + \epsilon_{1,t}
\end{equation}
The coefficient $c$ captures the total effect.

\textbf{Step 2:} Estimate the effect of geopolitical risk on ambiguity:
\begin{equation}
Amb_t = \alpha_2 + a \cdot GPRD_t + \gamma_2' X_t + \epsilon_{2,t}
\end{equation}
The coefficient $a$ measures how geopolitical risk influences the mediator.

\textbf{Step 3:} Estimate the effect of ambiguity on returns, controlling for geopolitical risk:
\begin{equation}
r_t = \alpha_3 + c' \cdot GPRD_t + b \cdot Amb_t + \gamma_3' X_t + \epsilon_{3,t}
\end{equation}
The coefficient $b$ measures the direct effect of ambiguity on returns, while $c'$ is the direct effect of geopolitical risk conditional on ambiguity.

The indirect effect is $a \times b$, and the proportion mediated is $(a \times b)/c$. We conduct bootstrap inference with 1,000 replications to test the significance of the indirect effect.

Additionally, we examine whether the impact of ambiguity on returns is moderated by market conditions. We estimate:
\begin{equation}
r_t = \alpha + \beta_1 Amb_t + \beta_2 Amb_t \times HighVol_t + \beta_3 Amb_t \times LowSent_t + \gamma' X_t + \epsilon_t
\end{equation}
where $HighVol_t$ is a dummy variable indicating the top quintile of volatility, and $LowSent_t$ indicates the bottom quintile of investor sentiment. Positive and significant coefficients on the interaction terms would suggest that ambiguity's negative effect on returns is amplified during stressed market conditions.

\begin{table}[htbp]
\centering
\caption{Mediation Analysis Results}
\label{tab:mediation_results}
\begin{tabular}{lcc}
\toprule
 & \multicolumn{1}{c}{Investor Sentiment} & \multicolumn{1}{c}{Market Uncertainty} \\ % Explicit column labels
\midrule
Path a: GPR → Mediator & $-0.142^{***}$ & $0.125^{***}$ \\
 & (-4.21) & (3.68) \\
Path b: Mediator → Returns & $0.049^{***}$ & $-0.058^{***}$ \\
 & (3.87) & (-4.52) \\
Direct Effect (c') & $-0.013$ & $-0.009$ \\
 & (-1.56) & (-1.08) \\
Indirect Effect (a×b) & $-0.007^{**}$ & $-0.007^{**}$ \\
 & [Sobel p = 0.032] & [Sobel p = 0.028] \\
Total Effect (c) & $-0.020^{**}$ & $-0.016^{*}$ \\
 & (-2.15) & (-1.72) \\
\midrule
Proportion Mediated & 35.0\% & 43.8\% \\
Observations & 5,832 & 5,832 \\
\bottomrule
\end{tabular}

\vspace{1em} % Adds a small vertical space (adjust 0.5em as needed, e.g., 1em for more space)
\footnotesize
\raggedright % Aligns the footnote to the left
Note: t-statistics in parentheses; Sobel test p-values in brackets. 
$^{***}$ p$<$0.01, $^{**}$ p$<$0.05, $^{*}$ p$<$0.1. 
Mediators are adapted for the Chinese stock market: Investor Sentiment (proxied by CSI 300 sentiment index) and Market Uncertainty (proxied by SSE 50 volatility index).
\end{table}

\subsubsection{Robustness Checks}

To ensure the robustness of our findings, we conduct several additional analyses:

\textbf{Alternative Ambiguity Measures.} We construct alternative proxies for ambiguity, including analyst forecast dispersion for Chinese firms, disagreement in professional forecasters' macroeconomic predictions, and volatility-of-volatility measures. Consistency across measures would strengthen confidence in our results.

\textbf{Subsample Analysis.} We split the sample into periods of high and low geopolitical risk to examine whether the relationships are stable or time-varying. We also examine subperiods corresponding to specific geopolitical episodes (e.g., 2018-2019 trade war, 2022 Russia-Ukraine conflict).

\textbf{Granger Causality.} We implement vector autoregression (VAR) models and Granger causality tests to examine the dynamic relationships and directional causality between geopolitical risk, ambiguity, and returns. Impulse response functions will trace out how shocks to geopolitical risk propagate through ambiguity to affect returns over multiple horizons.

\textbf{Event Studies.} We conduct event studies around specific, high-profile geopolitical events to examine the immediate and delayed responses of ambiguity and returns, providing micro-level evidence to complement our time-series analysis.

\textbf{Structural Break Tests.} We employ Chow tests and recursive parameter estimation to detect potential structural breaks in the relationships, ensuring that our findings are not driven by a specific sub-period.

\section{Conclusion}

This study provides evidence on how geopolitical risk transmits to the Chinese capital market through the channel of financial ambiguity. We construct a daily measure of ambiguity and show that it significantly explains stock returns, independent of traditional risk measures. Our findings reveal that geopolitical risk increases market ambiguity, which in turn negatively affects returns.

The mediation analysis shows that investor sentiment is a key mechanism through which geopolitical uncertainty affects prices. The moderation results highlight that institutional factors matter—state-owned enterprises are less vulnerable to ambiguity than private firms. These findings have important implications for investors needing to look beyond volatility when assessing market uncertainty, and for policymakers concerned about financial stability during geopolitical crises.

Future research could extend this work in several directions. Examining sector-specific effects—particularly comparing China's technology and manufacturing sectors—would provide insights into how different industries respond to geopolitical uncertainty. Additionally, exploring cross-country comparisons with other emerging markets would help determine whether our findings generalize to other contexts. Finally, developing higher-frequency ambiguity measures could improve our understanding of how quickly geopolitical information is incorporated into prices.

\bibliographystyle{plain}
\bibliography{geo}

\end{document}